\documentclass{book}
\usepackage{lipsum}

%%%%%%%%%%%%%%%%%%%%%%%%%%%%%%%%%%%%%%%%%%%%%%%%%%%%%%%%%%%%%%%%%%%%%%%%%%%%%%%%%%%%%%%%
% RPG BUNDLE PACKAGE
%%%%%%%%%%%%%%%%%%%%%%%%%%%%%%%%%%%%%%%%%%%%%%%%%%%%%%%%%%%%%%%%%%%%%%%%%%%%%%%%%%%%%%%%

% The bundle accepts explicit key/value options and routes them to its modules.
\usepackage[
    theme=maritime,
    alignment=raggedright,
    columns=two,
]{rpg}

%%%%%%%%%%%%%%%%%%%%%%%%%%%%%%%%%%%%%%%%%%%%%%%%%%%%%%%%%%%%%%%%%%%%%%%%%%%%%%%%%%%%%%%%

\begin{document}

\chapter{RPG Bundle Example}

\section{The Old Road}

The road follows a cold river through pine woods.  The \texttt{rpg} package loads the color, font, and layout modules together.  This example uses the maritime color theme, ragged-right body text, RPG heading styles, and a two-column page layout.

\lipsum[1]

\subsection{Travel Notes}

\begin{itemize}
    \item Blue lanterns mark the path to a ruined watchtower.
    \item The stones are carved with warnings in an older tongue.
\end{itemize}

\section{At the Tower}

The bundle interface intentionally uses key/value options. Shortcut options such
as \texttt{twocolumn} are available on \texttt{rpglayout}, but not on
\texttt{rpg}.

\lipsum[2]

\end{document}
